\section{Desarrollo del proyecto}

% =====================================================================
\subsection{Fase 0: preparación del entorno}

\subsubsection{Motor, módulos y plantilla}

El proyecto se creó en Unity 6.6 (versión 6000.6.3f1) con la plantilla \textit{Universal 3D}, que
usa el Universal Render Pipeline (URP). URP es el sistema de renderizado recomendado para proyectos
que se exportan a móviles y navegadores, porque permite iluminación, materiales con emisión y
posprocesado con un costo controlado. Además del soporte para Windows y Web, se instalaron los
módulos de Android: \textit{Android Build Support}, \textit{Android SDK \& NDK Tools} y
\textit{OpenJDK}, necesarios para generar el APK. Al proyecto se añadieron los paquetes de la
Tabla~\ref{tab:paquetes}.

\begin{TablaAPA}{tab:paquetes}{Paquetes de Unity utilizados}
\begin{tabularx}{\linewidth}{P{4.2cm}L}
\toprule
Paquete & Uso en Nexus \\
\midrule
Universal RP 17.6 & Renderizado, luces, materiales con emisión y posprocesado \\
Input System 1.19 & Teclado, ratón, mando y entradas táctiles \\
AI Navigation 2.0 & Mallas de navegación (NavMesh) y agentes para los androides \\
uGUI y TextMeshPro & Interfaz: menú, HUD, inventario y pantallas finales; texto en el espacio 3D \\
Pipeline & Control del editor desde la línea de comandos para automatizar pruebas y compilaciones \\
\bottomrule
\end{tabularx}
\end{TablaAPA}

\subsubsection{Organización del proyecto y escenas}

Todo el contenido propio se agrupó en la carpeta \cod{Assets/Nexus}, con una subcarpeta por tipo
de recurso: escenas, scripts, herramientas de editor, prefabs, materiales, texturas, modelos,
animaciones, audio, interfaz, datos y perfiles de compilación. Se crearon cuatro escenas, que se
cargan en este orden: \cod{00\_Menu} (menú principal), \cod{01\_Mantenimiento},
\cod{02\_Almacenamiento} y \cod{03\_Nucleo} (Figura~\ref{fig:estructura}).

\Figura[0.9\linewidth]{003_estructura-carpetas-y-escenas.jpg}{fig:estructura}{Estructura de carpetas del proyecto y escenas del juego}{}

\subsubsection{Configuración del reproductor y perfiles de compilación}

En \ruta{Project Settings} \textrightarrow{} \ruta{Player} se configuraron el nombre del producto
(\textit{Nexus}) y de la empresa (\textit{Nexus Team}), y los ajustes propios de cada plataforma
(Figura~\ref{fig:player}): en Windows, ventana redimensionable; en Android, el identificador
\cod{com.nexusteam.nexus}, el compilador IL2CPP con arquitectura ARM64, Android~8.0 como versión
mínima y orientación solo horizontal; en Web, compresión con descompresión de respaldo, para que el
juego funcione en cualquier servidor. Android y Web usan el nivel de calidad \textit{Mobile} y
Windows el nivel \textit{PC}.

\Figura[0.85\linewidth]{004_player-settings-android.jpg}{fig:player}{Ajustes del reproductor para Android}{}

Para generar las tres versiones desde el mismo proyecto se crearon tres perfiles de compilación
(\textit{Build Profiles}): \cod{Nexus\_PC\_Windows}, \cod{Nexus\_Web} y \cod{Nexus\_Android}
(Figura~\ref{fig:perfiles}). Cada perfil guarda la configuración de su plataforma, y cambiar de
plataforma consiste en activar el perfil correspondiente.

\Figura[0.85\linewidth]{005_build-profiles-plataformas.jpg}{fig:perfiles}{Ventana de perfiles de compilación con las tres plataformas}{}

% =====================================================================
\subsection{Fase 1: modelado de objetos y del entorno}

\subsubsection{Herramienta de construcción}

Los objetos del juego se modelaron combinando primitivas de Unity (cubos, cilindros y planos) con
texturas y materiales propios. Todo se genera con la herramienta de editor
\cod{ConstructorFase1}, que agrega al menú \ruta{Nexus} \textrightarrow{} \ruta{Fase 1} los pasos
\textit{Generar texturas}, \textit{Crear materiales} y \textit{Construir prefabs}. Cada objeto queda
guardado como prefab, listo para colocarlo en los niveles.

\subsubsection{Texturas procedurales}

Las texturas se generan por código, píxel a píxel, por lo que son originales. Junto con cada
textura de color se genera un mapa de alturas y, a partir de él, un mapa de normales que da relieve
a la luz: los agujeros se ven hundidos y las asas sobresalen (Tabla~\ref{tab:texturas}).

\begin{TablaAPA}{tab:texturas}{Texturas generadas por código}
\begin{tabularx}{\linewidth}{P{3.6cm}P{2.3cm}L}
\toprule
Textura & Tamaño & Contenido \\
\midrule
Frente de rack & 256 × 1024 & Rieles con agujeros de montaje, diez servidores con malla perforada, etiquetas, huecos para LEDs, botones, asas y rejilla de ventilación \\
Piso técnico & 512 × 512 & Baldosas de 60~cm con tornillos; una de cada cuatro perforada para la ventilación \\
Pared & 512 × 512 & Paneles metálicos con juntas, remaches, banda horizontal y zócalo \\
Franjas de peligro & 256 × 256 & Diagonales amarillas y negras para los marcos de las puertas \\
Metal cepillado & 256 × 256 & Vetas horizontales para el cuerpo de los racks \\
Tarjetas de acceso & 256 × 160 & Banda, chip dorado y logotipo, en rojo y en azul \\
\bottomrule
\end{tabularx}
\end{TablaAPA}

\subsubsection{Sombreador de LEDs}

Para los cientos de luces de los racks se escribió un sombreador propio en HLSL para URP. El
parpadeo se calcula en la tarjeta gráfica a partir del tiempo, sin ningún script, lo que mantiene
el juego fluido en Android y WebGL. Cada LED guarda en el color de sus vértices su color (verde,
azul, ámbar o rojo) y un desfase, de modo que todos los LEDs de un rack forman una sola malla y aun
así cada uno parpadea a su ritmo. El color se multiplica por una intensidad alta para que el efecto
\textit{Bloom} del posprocesado los haga brillar. El mismo sombreador, en modo de pulso suave, se
usa en el neón de los discos, en los lectores de las puertas y en las pantallas de las terminales.

\subsubsection{Objetos del centro de datos}

\begin{itemize}
  \item \textbf{Rack de servidores} (0,6 × 2,0 × 1,0~m): cuerpo de metal cepillado, frente y
        trasera con la textura de rejilla y una malla con 80 LEDs colocados exactamente sobre los
        huecos de la textura (Figura~\ref{fig:rack}, izquierda). Los racks se agrupan en filas de diez,
        combinadas en tres mallas por celda, con lo que un laberinto de cientos de racks se dibuja
        con pocas llamadas a la tarjeta gráfica.
  \item \textbf{Disco Duro Maestro}: cubo metálico con un anillo y placas de neón azul y una luz
        propia; gira y flota, y se recoge al tocarlo (Figura~\ref{fig:rack}, derecha).
  \item \textbf{Tarjetas de acceso} roja y azul, que abren las puertas de su color.
  \item \textbf{Puerta de seguridad}: marco con franjas de peligro, dos hojas que se deslizan y un
        lector cuya luz indica qué tarjeta se necesita (Figura~\ref{fig:puerta}, izquierda). Hay variantes
        roja, azul, libre y sin energía.
  \item \textbf{Terminal}: consola con teclado y monitor con texto real en la pantalla, escrito con
        TextMeshPro en el espacio 3D (Figura~\ref{fig:puerta}, derecha).
  \item \textbf{Módulos de 3 × 3~m} de piso y pared, con los que se arman todos los niveles.
  \item \textbf{Objetos del nivel}: fusible, memoria USB, bebida energética, barrita de proteína,
        cartucho de refrigerante, mesa con nota, panel de fusibles, terminal de respaldo, válvula de
        refrigeración y ascensor. Los objetos que se recogen tienen una leve emisión para distinguirse
        en la oscuridad sin gastar una luz cada uno.
\end{itemize}

\FiguraDoble{014_rack-servidor-detalle.jpg}{015_disco-duro-maestro.jpg}{fig:rack}{Racks de servidores con sus LEDs (izquierda) y Disco Duro Maestro sobre su pedestal (derecha)}{}

\FiguraDoble{016_puerta-seguridad-roja.jpg}{017_terminal-nexus.jpg}{fig:puerta}{Puerta de seguridad roja (izquierda) y terminal con texto en pantalla (derecha)}{}

Los productos de curación tienen un diseño propio, sin marcas reales: una lata de bebida energética
y una barrita de proteína con envoltura de grafito, bordes plateados y una franja turquesa
(Figura~\ref{fig:curacion}).

\Figura[0.6\linewidth]{111_bebida-energetica-y-barrita.jpg}{fig:curacion}{Bebida energética y barrita de proteína}{}

\subsubsection{Utilería y decoración automática}

Para que los escenarios se vean como una empresa abandonada a la fuerza se modelaron 48 objetos de
utilería, también con primitivas y texturas generadas por código (Figura~\ref{fig:utileria}):
escritorios con computadores, sillas, archivadores, casilleros, una consola de cámaras de seguridad,
pizarras y carteles de la empresa; máquinas expendedoras, mesas y una nevera para la cafetería;
estanterías, cajas, bobinas de cable y gabinetes eléctricos para el almacén; y papeles, latas,
bolsas de basura, paneles del techo caídos, manchas, charcos y conos de precaución. Se agregaron
efectos de polvo flotante, chispas y fluorescentes que titilan. Las piezas de cada objeto se
combinan en una sola malla por material para que su costo de dibujo sea bajo.

\FiguraDoble{056_utileria-vitrina-general.jpg}{057_utileria-maquinas-y-escritorios.jpg}{fig:utileria}{Utilería del centro de datos: vitrina general (izquierda) y máquinas expendedoras con escritorios (derecha)}{}

Un decorador automático amuebla cada sala según su tema (oficina, seguridad, cafetería, almacén,
servidores, eléctrico, mantenimiento o pasillo). Primero coloca los muebles obligatorios del tema,
luego llena las paredes que ve la cámara de la sala hasta una densidad definida, cuelga carteles y
gabinetes y, por último, riega basura por el suelo y añade polvo del tamaño de la sala. Los muebles
van siempre contra las paredes, para no estorbar el paso del jugador ni las rutas de los androides,
y nunca junto a una cámara o una puerta. Los muebles altos bloquean la visión de los androides, así
que también sirven para esconderse. La decoración usa una semilla fija: reconstruir un nivel da
siempre el mismo resultado.

Las máquinas expendedoras dan bebidas energéticas y barritas de proteína hasta agotarse, pero el
producto cae con un golpe que atrae a los androides que estén a menos de 11~m, lo que convierte la
curación en una decisión con riesgo.

% =====================================================================
\subsection{Fase 2: personajes y animación}

\subsubsection{Avatares de los integrantes}

Los dos personajes jugables son los avatares de los integrantes del equipo, creados con Avaturn a
partir de fotografías \parencite{avaturn2026}. Se descargaron en formato GLB, que guarda en un solo
archivo la malla, un esqueleto de 52 huesos con nombres estándar y las texturas. Con un script de
Python para Blender \parencite{blender2026} se convirtieron a FBX con las texturas incrustadas, sin
huesos adicionales y con la escala en metros, que es el formato con el que Unity configura los
avatares humanoides.

\subsubsection{Configuración como Avatar Humanoid}

En Unity, los modelos se importaron con el tipo de animación \textit{Humanoid} y la definición
\textit{Create From This Model} (Figura~\ref{fig:rig}). Unity generó un avatar válido para cada
integrante, con los 52 huesos asignados al esqueleto humano estándar (Figura~\ref{fig:avatares}).
El sistema Humanoid traduce el esqueleto de cada modelo a un esqueleto común; gracias a eso, una
misma animación funciona en los dos personajes y en los androides aunque tengan proporciones
distintas, y es lo que permite que el selector de personajes use un único controlador de animación
\parencite{unitymanual2026}.

\Figura[0.75\linewidth]{007_rig-humanoid-neyder.jpg}{fig:rig}{Pestaña Rig del modelo de Neyder con el tipo de animación Humanoid}{}

\FiguraDoble{008_configure-avatar-mapeo-huesos-neyder.jpg}{009_configure-avatar-mapeo-huesos-deisy.jpg}{fig:avatares}{Mapeo de huesos del Avatar Humanoid de Neyder (izquierda) y de Deisy (derecha)}{}

Las texturas se asignaron a materiales URP con mapa de color y de normales. El cabello usa recorte
por transparencia (\textit{alpha clipping}) y se dibuja por las dos caras, una solución más ligera
que la transparencia real para móviles y navegadores; la piel y la ropa usan un flujo de trabajo no
metálico (Figura~\ref{fig:integrantes}).

\Figura[0.85\linewidth]{006_avatares-integrantes-en-unity.jpg}{fig:integrantes}{Avatares de los integrantes en Unity}{}

\subsubsection{Animaciones y controlador}

Los integrantes descargaron de Mixamo \parencite{mixamo2026} las animaciones aplicadas sobre sus
avatares. Se eligieron once clips: reposo, caminar, correr, caminar herido, reacción a un golpe,
golpe al caminar, morir, y reposo, caminar, disparar y apuñalar con arma. Cada clip se configuró
como Humanoid; el avance de la cadera se deja fuera del personaje (\textit{root motion}
descartado), de modo que la animación se reproduce en su sitio y es el componente
\textit{CharacterController} el que mueve al técnico.

El controlador de animación \cod{AC\_Tecnico}, compartido por los dos personajes, tiene un árbol de
mezcla de locomoción (reposo, caminar a 1,8~m/s y correr a 5~m/s) cuyos umbrales coinciden con la
velocidad real de las animaciones, para que los pies no se deslicen sobre el suelo; una variante
herida; los estados de golpe y muerte; y los estados y parámetros de las armas. Al retroceder, la
animación de caminar se reproduce en reversa (Figura~\ref{fig:locomocion}).

\FiguraDoble{022_jugador-animacion-reposo-modo-play.jpg}{023_jugador-caminando-hacia-camara-fija.jpg}{fig:locomocion}{Animación de reposo (izquierda) y el técnico caminando hacia la cámara fija (derecha)}{}

\subsubsection{Personajes secundarios y cadáveres}

De Mixamo se obtuvieron también dos operadores tácticos y un empleado con casco y chaleco, que en
el juego aparecen como los cuerpos de quienes entraron antes que el técnico. Un script de Blender
reduce sus texturas a un máximo de 1024~píxeles, separa los canales de color, normales, rugosidad y
emisión, y combina la rugosidad en el formato que usa el material URP; con esto los archivos pasaron
de hasta 115~MB a unos 10~MB. En Unity se configuraron como Humanoid.

Las poses de los cadáveres se describen con los músculos del sistema Humanoid, valores de $-1$ a 1
que funcionan con cualquier modelo: una pose tendida, tomada del último cuadro de la animación de
muerte, y una pose sentada contra la pared, creada músculo por músculo (piernas al frente, una
rodilla doblada, torso caído y cabeza ladeada). La pose se escribe en los huesos, de modo que el
prefab no necesita un controlador de animación. Cada cuerpo tiene un charco de sangre, una
colisión baja que los androides rodean y una luz tenue para distinguirse en la oscuridad
(Figura~\ref{fig:cadaveres}).

\FiguraDoble{082_cadaveres-tendidos.jpg}{083_cadaveres-sentados-pared.jpg}{fig:cadaveres}{Poses de cadáver tendidas (izquierda) y sentadas contra la pared (derecha), creadas con los músculos del avatar}{}

\subsubsection{Animaciones de la escopeta creadas por código}

Para la escopeta de corredera se crearon animaciones propias a partir de la técnica de manejo real
de este tipo de arma: culata en el hombro, mano fuerte en la empuñadura y mano de apoyo en el
guardamanos; posición de listo con el cañón hacia abajo; retroceso hacia el hombro; bombeo para
cargar el siguiente cartucho; y recarga a la altura de la cadera, empujando los cartuchos uno por
uno.

\begin{enumerate}
  \item Se escribieron por código, músculo por músculo, tres clips Humanoid: listo, apuntar y
        recargar.
  \item Se agregó al controlador una capa \textit{Escopeta} con una máscara de brazos, manos y
        cabeza: las piernas siguen caminando y corriendo con la capa base.
  \item La cinemática inversa (IK) coloca el arma respecto al hombro según la pose y lleva las
        manos a la empuñadura y al guardamanos. El retroceso, el bombeo y el ciclo de cada cartucho
        se calculan con curvas en el código.
  \item La recarga cuesta tiempo real: 0,5~s por cartucho.
\end{enumerate}

\FiguraDoble{088_escopeta-listo.jpg}{089_escopeta-apuntar-derecha.jpg}{fig:escopeta}{Escopeta en posición de listo (izquierda) y apuntando con la mano de apoyo en el guardamanos (derecha)}{}

\Figura[0.7\linewidth]{091_escopeta-recarga.jpg}{fig:escopeta-recarga}{Recarga de la escopeta, cartucho por cartucho}{}
